\section{Integration Approach}
\label{sec:integration_approach}%

The integration testing strategy for the Best Bike Paths system follows a \textbf{Bottom-Up approach}, as defined in the Design Document (see \Cref{ch:implementation_integration_test_plan} of the DD). This approach ensures that foundational components are tested and validated before higher-level components are integrated.

The Bottom-Up strategy provides several advantages:
\begin{itemize}
    \item Early validation of core data and service layers
    \item Easier debugging by isolating issues at lower levels first
    \item Support for parallel development across different components
    \item Reduced integration effort by ensuring stable dependencies
\end{itemize}

\section{Integration Levels}
\label{sec:integration_levels}%

The integration process is organized into the following levels:

\subsection{Level 1: Data Layer Integration}
\label{subsec:level1_data_layer}%
Integration of database components, repositories, and data access layers. This includes:
\begin{itemize}
    \item DBMS (PostgreSQL) integration
    \item Repository abstractions
    \item Cache layer (Redis) integration
    \item Data migration and seed data validation
\end{itemize}

\subsection{Level 2: Service Layer Integration}
\label{subsec:level2_service_layer}%
Integration of individual microservices with the data layer and with each other. This includes:
\begin{itemize}
    \item Core services (User Management, Trip Management, Path Management)
    \item Advanced services (Sensor Data Processing, Path Search, Rating, Information Merging)
    \item Service-to-service communication
    \item Message queue integration (RabbitMQ)
\end{itemize}

\subsection{Level 3: API Gateway Integration}
\label{subsec:level3_api_gateway}%
Integration of the API Gateway with all backend services. This includes:
\begin{itemize}
    \item Request routing and versioning
    \item Authentication and authorization
    \item Request validation
    \item Error handling and response formatting
\end{itemize}

\subsection{Level 4: External Services Integration}
\label{subsec:level4_external_services}%
Integration with third-party services. This includes:
\begin{itemize}
    \item Weather API integration
    \item Map Service API integration
    \item Email Service (SMTP) integration
    \item Fallback and error handling mechanisms
\end{itemize}

\subsection{Level 5: Client Integration}
\label{subsec:level5_client_integration}%
Integration of the mobile application with backend services. This includes:
\begin{itemize}
    \item Mobile application core components
    \item Sensor Manager integration
    \item Trip Detection Engine integration
    \item Obstacle Detection Module integration
    \item Client-Server communication
\end{itemize}

\subsection{Level 6: End-to-End Integration}
\label{subsec:level6_e2e_integration}%
Complete system integration testing through use case scenarios. This includes:
\begin{itemize}
    \item Critical use cases (UC1-UC3, UC6)
    \item Advanced use cases (UC4-UC5, UC7-UC12)
    \item Error scenarios and edge cases
\end{itemize}

\section{Integration Order}
\label{sec:integration_order}%

The integration order follows the Bottom-Up approach and aligns with the implementation phases defined in the Design Document:

\begin{enumerate}
    \item \textbf{Phase 0:} Project foundations (repository structure, domain model, database schema)
    \item \textbf{Phase 1:} Core services integration (User Management, Trip Management, Path Management, Rating Service)
    \item \textbf{Phase 2:} API Gateway integration and API stabilization
    \item \textbf{Phase 3:} Client flows and discovery features (Path Search Service, Map integration, Caching)
    \item \textbf{Phase 4:} Automated detection and sensor processing (Mobile sensor layer, Sensor Data Processing Service)
    \item \textbf{Phase 5:} Information merging (Information Merging Service integration)
\end{enumerate}

\section{Testing Strategy}
\label{sec:testing_strategy}%

The testing strategy follows a pyramid approach, supporting the Bottom-Up integration order:

\subsection{Unit Tests}
\label{subsec:unit_tests}%
Unit tests validate individual components in isolation before integration. These tests focus on:
\begin{itemize}
    \item Business logic within services
    \item Repository operations using in-memory fakes
    \item Component-specific functionality
\end{itemize}

\subsection{Integration Tests}
\label{subsec:integration_tests}%
Integration tests validate component interactions at each integration level:
\begin{itemize}
    \item Service-to-DB integration tests
    \item Service-to-Service integration tests
    \item Gateway-to-Service integration tests
    \item External service integration tests
    \item Client-Server integration tests
\end{itemize}

\subsection{End-to-End Tests}
\label{subsec:e2e_tests}%
End-to-end tests validate complete user flows through the system:
\begin{itemize}
    \item Critical use case flows (UC1-UC3, UC6)
    \item Advanced use case flows (UC4-UC5, UC7-UC12)
    \item Error handling scenarios
    \item Offline/online synchronization
\end{itemize}

\subsection{Non-Functional Tests}
\label{subsec:non_functional_tests}%
Non-functional tests validate system quality attributes:
\begin{itemize}
    \item Performance smoke tests (response times, cache hit rates)
    \item Security checks (authentication, authorization, input validation)
    \item Resilience checks (graceful degradation, error recovery)
\end{itemize}
